\documentclass{article}
\usepackage[utf8]{inputenc}
\usepackage{amsfonts}
\usepackage{amsmath}

\title{Schr\"odinger equation}
\author{Jean-Marc Nuzillard}
\date{May 2026}

\begin{document}

\maketitle

\section{Introduction}
Assuming a matter particle may be described by a wave (a Louis de Broglie wave), this wave must 
obey to a wave equation, like sound, light, or heat waves obey to an equation.
The path along which the Schr\"{o}dinger equation (SE) for matter waves
was obtained is somehow tortuous and no clear historical account of it is available.
The SE may considered as an {\em axiom} and therefore has no need for a justification,
apart the fact that it works when it must work.
As an undergraduate student, our physics teacher attempted nevertheless to give us an idea of how the SE can be derived,
even though this derivation must not be taken as a {\em proof}.

The derivation of the SE starts with an assumption: a particle of mass $m$ is described by
a function (the wave function) of the three spatial coordinates $x$, $y$, $z$
and of the time coordinate. This function, named $\psi$ has complex values.
Otherwise stated, the wave function is a complex-valued scalar field.

\[
\psi(x, y, z, t) \in \mathbb{C}
\]

When dealing with waves, everything becomes simpler with complex numbers.
A particle, something that one might think of as a mathematical point 
is described by a function that takes values in the whole three-dimensional space, at each time.
The values of this function evolve during time, moreover in a deterministic way.
So, if $\psi$ is known at some time $t$, then the SE should provide a way 
to know what it becomes later.

Knowing the value of $\psi$ at some time gives access to the density of presence probability of the particle.
Considering a small (infinitesimal) volume $\mathrm{d}V$ around the point of coordinates $(x, y, z)$ then
the probability of presence of the particle is

\[
\mathrm{d}p = |\psi(x,y,z,t)|^2 \mathrm{d}V
\]

The particle is therefore everywhere in space at any time but not everywhere with the same probability.
The classical trajectory is replaced by the time evolution of the probability density function.

The simpliest possible wavefunction is the planar sine wave. 
A planar wave is defined by a direction of propagation, hereafter considered as the $x$-axis 
of a set of orthogonal axis that defines space cartesian coordinates.
This trick reduces what follows to what would happen in a 1-dimensional space.
The SE is first established here for a planar wave, assuming that any physically meaningful wave function 
can be decomposed into a sum of planar wave functions.
This decomposition is a kind of generalized Fourier transformation and holds because
we expect the SE to be a linear differential equation, what it is.

So, we assume that an elementary sine planar, complex-valued, wave-function may be written as

\[
\psi(x, y, z, t) = \exp\left( 2\pi i \left[ \frac{x}{\lambda} - \frac{t}{T} \right] \right)
\]

\noindent in which the spatial periodicity of the wave (or wavelength) is $\lambda$
and the temporal periodicity (or period) is $T$.
It can be checked that $\psi(x+\lambda, t)=\psi(x, t)$ and $\psi(x, t+T)=\psi(x,t)$.
The minus sign in the expression of the elementary wave function indicates a propagation along the $+x$ axis.

How can be related the wave parameters $\lambda$ and $T$ to the mechanical behavior of a particle of mass $m$?

Louis de Broglie associated a particle of momentum $p$ (product of mass and speed)
with a matter wave of wavelength $\lambda$ so that

\[
\lambda = \frac{h}{p}
\]

\noindent in which $h$ is the Planck constant.
The planet-model of the hydrogen atom in which the electron is submitted to the constraint of
having the length of electron orbit being a integer multiple of the associated wavelength 
provides the correct stability energy levels, as spectroscopically proved by measured differences between energy levels.

The time periodicity defined by the period $T$ is related to the frequency $\nu = 1/T$. 
The frequency of light wave is related to the energy quantum it carries by the Planck's relationship.
It is assumed here that this relation holds identically for matter waves.

\[
E = h\nu = \frac{h}{T} \quad \mbox{or} \quad T = \frac{h}{E}
\]

\noindent in which E is now the total mechanical energy of the particle.

The expression of the planar $\psi$ function may be rewritten

\[
\psi(x, y, z, t) = \exp\left(i \left[ x \frac{2\pi}{\lambda} - t \frac{2\pi}{T} \right] \right)
\]

\noindent that reduces to

\[
\psi(x, y, z, t) = \exp\left(i \left[ k x - \omega t \right] \right) \quad \mbox{with} \quad k=\frac{2\pi}{\lambda} \quad \mbox{and} \quad \omega=\frac{2\pi}{T}
\]

With the reduced Planck constant $\hbar = h/2\pi$

\[
E = \hbar k \quad \mbox{and} \quad p = \hbar k
\]

The latter equations associate the matter wave parameters $k$ and $\omega$, the spatial and temporal pulsations of a planar wave,
with mechanical parameters of the corresponding particle, impulsion and energy.

Energy here is the total mechanical energy.
A new hypothesis is that the particle under consideration is submitted to forces
that derive from a potential energy.
Therefore its total energy is the sum of its kinetic energy $T$ and of a potentiel energy $V$ so that

\[
E = T + V \quad \mbox{or} \quad E = \frac{p^2}{2m} + V
\]

The law of the conservation of the total energy states that total energy of an isolated system is constant while
kinetic and potential energies may interconvert.

\section{Assembling the pieces}
The derivative of $\psi$ as a function of the spatial coordinate $x$ gives

\[ 
\frac{\partial \psi}{\partial x} = ik\psi
\]

\noindent so that

\[
-i \hbar \frac{\partial \psi}{\partial x} = p \psi
\]

The particle impulsion appears as a multiplicative factor to $\psi$ by application of the operator

\[
\hat{p} = -i \hbar \frac{\partial}{\partial x}
\]

\noindent on $\psi$ itself.
This does not directly give access to $p$ but to $p\psi$.
Accessing the kinetic energy $T$ is possible through $p^2 \psi$, applying twice the operator associated to $p$:

\[
- \hbar^2 \frac{\partial^2 \psi}{\partial x^2} = p^2 \psi
\]

\noindent and

\[
- \frac{\hbar^2}{2m} \frac{\partial^2 \psi}{\partial x^2} = T \psi
\]

Pushing further the same idea, the derivative of $\psi$ as a function of time gives

\[ 
\frac{\partial \psi}{\partial t} = -i \omega \psi
\]

\noindent so that

\[
i \hbar \frac{\partial \psi}{\partial t} = E \psi
\]

Multipling by $\psi$ the two terms of the classical expression of the total energy yields:

\[
T\psi + V\psi = E\psi
\]

\indent or

\[
- \frac{\hbar^2}{2m} \frac{\partial^2 \psi}{\partial x^2} + V\psi = i \hbar \frac{\partial \psi}{\partial t}
\]

This equation is the one-dimensional version of the SE.

In the general case, there are 3 space coordinates $x$, $y$, and $z$, and the particle momentum is a vector 
with 3 coordinates $p_x$, $p_y$, and $p_z$, each associated to an operator:

\[
\hat{p_x}=-i \hbar \frac{\partial}{\partial x} \quad \hat{p_y}=-i \hbar \frac{\partial}{\partial y} \quad \hat{p_z}=-i \hbar \frac{\partial}{\partial z}
\]

The action of the operator associated to $p^2 = p_x^2 + p_y^2 + p_z^2$ can be written for any planar matter wave

\[
- \hbar^2  \left( \frac{\partial^2 \psi}{\partial x^2} + \frac{\partial^2 \psi}{\partial y^2} + \frac{\partial^2 \psi}{\partial z^2} \right) =  p^2 \psi 
\]

\noindent also written

\[
- \hbar^2 \Delta \psi = p^2 \psi
\]

\noindent by the introduction of the Laplacian operator $\Delta$ that sums the three second partial derivatives of a function of space coordinates.

Therefore, any planar matter wave, any superposition of planar waves, 
and therefore any wave function associated to a matter particle obey the Schr\"{o}dinger equation:

\[
\boxed{
- \frac{\hbar^2}{2m} \Delta \psi + V\psi = i \hbar \frac{\partial \psi}{\partial t}
}
\]
\end{document}
